\documentclass[12pt,openany,oneside,a4paper,english,brazil]{abntex2}

% =============================================================
% DADOS PRINCIPAIS DO TCC
% Edite este arquivo para atualizar capa, folha de rosto e banca.
% =============================================================

\newcommand{\autorTCC}{Fabrício Miguel Novicki}
\newcommand{\tituloTCC}{Desenvolvimento e validação de um algoritmo para otimização geométrica de \textit{propellers} com base em critérios estruturais, aerodinâmicos e aeroacústicos}
\newcommand{\cidadeTCC}{Florianópolis}
\newcommand{\anoTCC}{2027} % Ajustar caso a defesa/depósito ocorra em outro ano.

\newcommand{\universidadeTCC}{Universidade Federal de Santa Catarina}
\newcommand{\centroTCC}{Centro Tecnológico}
\newcommand{\departamentoTCC}{Departamento de Engenharia Mecânica}
\newcommand{\cursoTCC}{Curso de Engenharia Mecânica}

\newcommand{\orientadorTCC}{Prof. Júlio Apolinário Cordiolli, Dr. Eng.}
\newcommand{\coorientadorTCC}{Lucas Meirelles Pereira, M. Eng.}

% Dados a completar perto da defesa
\newcommand{\dataDefesaTCC}{[dia] de [mês] de [ano]}
\newcommand{\coordenadorTCC}{Prof. [nome do coordenador], Dr. Eng.}
\newcommand{\membroBancaUm}{Prof. [nome do membro], Dr. Eng.}
\newcommand{\instituicaoBancaUm}{[instituição]}
\newcommand{\membroBancaDois}{Prof. [nome do membro], Dr. Eng.}
\newcommand{\instituicaoBancaDois}{[instituição]}

% =============================================================
% CONFIGURAÇÃO GERAL - UFSC / ABNT
% Compilar com pdfLaTeX + BibTeX (ou latexmk -pdf).
% =============================================================

\usepackage[T1]{fontenc}
\usepackage[utf8]{inputenc}
% O template Word utiliza Arial. Helvetica (pacote helvet) é a aproximação
% livre e portátil adotada aqui para compilar sem arquivos de fonte externos.
\usepackage[scaled=0.97]{helvet}
\renewcommand{\familydefault}{\sfdefault}

\usepackage{microtype}
\usepackage{graphicx}
\usepackage{booktabs}
\usepackage{longtable}
\usepackage{tabularx}
\usepackage{array}
\usepackage{multirow}
\usepackage{amsmath,amssymb,mathtools,bm}
\usepackage{siunitx}
\usepackage{float}
\usepackage{enumitem}
\usepackage{caption}
\usepackage{subcaption}
\usepackage{xcolor}
\usepackage{url}
\usepackage{hyperref}
\usepackage[alf]{abntex2cite}

% -------------------------------------------------------------
% Página e parágrafos - conforme o template fornecido
% Superior/esquerda: 3 cm; inferior/direita: 2 cm.
% Corpo: 12 pt, espaçamento 1,5; recuo de primeira linha: 1,5 cm.
% -------------------------------------------------------------
\setlrmarginsandblock{3cm}{2cm}{*}
\setulmarginsandblock{3cm}{2cm}{*}
\checkandfixthelayout

\setlength{\parindent}{1.5cm}
\setlength{\parskip}{0pt}
\OnehalfSpacing

% -------------------------------------------------------------
% Títulos
% -------------------------------------------------------------
\renewcommand{\ABNTEXchapterfont}{\bfseries\normalsize}
\renewcommand{\ABNTEXchapterfontsize}{\normalsize}
\renewcommand{\ABNTEXsectionfont}{\normalfont\normalsize}
\renewcommand{\ABNTEXsectionfontsize}{\normalsize}
\renewcommand{\ABNTEXsubsectionfont}{\normalfont\normalsize}
\renewcommand{\ABNTEXsubsectionfontsize}{\normalsize}
\renewcommand{\ABNTEXsubsubsectionfont}{\normalfont\normalsize}
\renewcommand{\ABNTEXsubsubsectionfontsize}{\normalsize}

% Capítulos iniciam em nova página (comportamento padrão da classe).
\setbeforesecskip{1.5\onelineskip}
\setaftersecskip{0.5\onelineskip}


% -------------------------------------------------------------
% Aparência dos capítulos e dos títulos pré-textuais.
%
% Capítulos numerados: alinhados à esquerda, no formato "1 - INTRODUÇÃO".
% Elementos pré-textuais não numerados: centralizados, em negrito e caixa alta,
% conforme o template da Engenharia Mecânica/UFSC.
% -------------------------------------------------------------
\renewcommand{\printchaptername}{}
\renewcommand{\chapternamenum}{}
\renewcommand{\printchapternum}{\thechapter.\space}
\renewcommand{\afterchapternum}{}

\makeatletter
\renewcommand{\printchaptertitle}[1]{%
  \ifthenelse{\boolean{abntex@innonumchapter}}{%
    \centering\bfseries\normalsize\MakeUppercase{#1}%
  }{%
    \raggedright\bfseries\normalsize\MakeUppercase{#1}%
  }%
}
\makeatother

% O ambiente de resumo do abnTeX2 usa uma formatação própria para o título.
% A definição abaixo garante o mesmo padrão: negrito, centralizado e caixa alta.
\renewcommand{\abstractnamefont}{\bfseries\normalsize}

% -------------------------------------------------------------
% Legendas
% Figuras: legenda abaixo (controle pela posição de \caption).
% Tabelas: legenda acima (controle pela posição de \caption).
% -------------------------------------------------------------
\captionsetup{font=small,labelfont=normalfont,labelsep=endash,justification=centering}

% -------------------------------------------------------------
% Numeração no canto superior direito apenas a partir da Introdução.
% -------------------------------------------------------------
\makepagestyle{ufsc}
\makeoddhead{ufsc}{}{}{\normalsize\thepage}
\makeevenhead{ufsc}{}{}{\normalsize\thepage}
\makeoddfoot{ufsc}{}{}{}
\makeevenfoot{ufsc}{}{}{}
\aliaspagestyle{chapter}{ufsc}

% -------------------------------------------------------------
% Links discretos no PDF.
% -------------------------------------------------------------
\hypersetup{
    pdftitle={\tituloTCC},
    pdfauthor={\autorTCC},
    colorlinks=true,
    linkcolor=black,
    citecolor=black,
    urlcolor=black,
    bookmarksdepth=4
}

% -------------------------------------------------------------
% Atalhos úteis
% -------------------------------------------------------------
\newcommand{\propeller}{\textit{propeller}}
\newcommand{\propellers}{\textit{propellers}}
\newcommand{\todo}[1]{\textbf{[TODO: #1]}}


\begin{document}
\selectlanguage{brazil}
\frenchspacing

% Títulos dos elementos pré-textuais conforme o template UFSC.
\renewcommand{\agradecimentosname}{AGRADECIMENTOS}
\renewcommand{\resumoname}{RESUMO}
\renewcommand{\listfigurename}{LISTA DE FIGURAS}
\renewcommand{\listtablename}{LISTA DE TABELAS}
\renewcommand{\listadesiglasname}{LISTA DE ABREVIATURAS E SIGLAS}
\renewcommand{\listadesimbolosname}{LISTA DE SÍMBOLOS}
\renewcommand{\contentsname}{SUMÁRIO}

% =============================================================
% CAPA - não contada e não paginada.
% =============================================================
\pagenumbering{gobble}
\begin{capa}
\begin{center}
    \includegraphics[width=2.4cm]{figures/ufsc_logo.png}

    \vspace{0.5cm}
    {\normalsize\MakeUppercase{\universidadeTCC}\par}
    {\normalsize\MakeUppercase{\centroTCC}\par}
    {\normalsize\MakeUppercase{\departamentoTCC}\par}
    {\normalsize\MakeUppercase{\cursoTCC}\par}

    \vfill
    {\normalsize \MakeUppercase\autorTCC\par}

    \vfill
    {\bfseries\normalsize\MakeUppercase{\tituloTCC}\par}

    \vfill
    {\normalsize \cidadeTCC\par}
    {\normalsize \anoTCC\par}
\end{center}
\end{capa}

\clearpage

% A partir da folha de rosto, páginas pré-textuais são contadas,
% mas o número não é exibido.
\pagenumbering{arabic}
\setcounter{page}{1}
\pretextual
\pagestyle{empty}

% Folha de rosto ("contracapa") organizada de acordo com o layout utilizado
% no Planejamento de TCC do autor: nome centralizado na parte superior,
% natureza do trabalho e orientação no semiplano direito e local/ano ao final.
\begin{folhaderosto}

\begin{center}
    \vspace*{2.4cm}
    {\normalsize\MakeUppercase{\autorTCC}\par}
\end{center}

\vspace*{4.6cm}

\begin{flushright}
    \begin{minipage}{0.50\textwidth}
        \OnehalfSpacing
        Trabalho de Conclusão de Curso apresentado ao Curso de Graduação em Engenharia Mecânica da Universidade Federal de Santa Catarina como requisito parcial para a obtenção do título de Bacharel em Engenharia Mecânica.

        \vspace{0.8cm}
        Orientador: \orientadorTCC

        \vspace{0.45cm}
        Coorientador: \coorientadorTCC
    \end{minipage}
\end{flushright}

\vfill
\begin{center}
    {\normalsize \cidadeTCC\par}
    {\normalsize \anoTCC\par}
\end{center}

\end{folhaderosto}

\clearpage
\thispagestyle{empty}
\vspace*{\fill}
\begin{center}
\begin{minipage}{0.78\textwidth}
\centering
\textbf{Ficha catalográfica para trabalhos acadêmicos}

\vspace{0.8cm}
\textit{[Elemento obrigatório na versão final.]}

\vspace{0.5cm}
Insira neste espaço a ficha catalográfica gerada pela Biblioteca Universitária da UFSC.
\end{minipage}
\end{center}
\vspace*{\fill}
\clearpage

\clearpage
\thispagestyle{empty}
\begin{center}
    {\normalsize \autorTCC\par}

    \vspace{1.5cm}
    {\bfseries\normalsize\MakeUppercase{\tituloTCC}\par}

    \vspace{1.5cm}
    Este Trabalho de Conclusão de Curso foi julgado adequado para obtenção do título de Engenheiro Mecânico e aprovado em sua forma final pelo Curso de Engenharia Mecânica.

    \vspace{0.8cm}
    \cidadeTCC, \dataDefesaTCC.

    \vspace{1.4cm}
    \rule{9cm}{0.4pt}\\
    \coordenadorTCC\\
    Coordenador do Curso

    \vspace{1.2cm}
    \textbf{Banca examinadora}

    \vspace{1.2cm}
    \rule{9cm}{0.4pt}\\
    \orientadorTCC\\
    Orientador

    \vspace{1.2cm}
    \rule{9cm}{0.4pt}\\
    \membroBancaUm\\
    \instituicaoBancaUm

    \vspace{1.2cm}
    \rule{9cm}{0.4pt}\\
    \membroBancaDois\\
    \instituicaoBancaDois

    \vfill
    \cidadeTCC, \anoTCC
\end{center}
\clearpage


% Elementos opcionais: mantenha, edite ou comente as linhas abaixo.
\begin{dedicatoria}
\vspace*{\fill}
\begin{flushright}
\begin{minipage}{0.50\textwidth}
\raggedleft
\todo{Inserir dedicatória, se desejado.}
\end{minipage}
\end{flushright}
\end{dedicatoria}

\begin{agradecimentos}
\todo{Escrever os agradecimentos na etapa final do TCC.}
\end{agradecimentos}

\begin{epigrafe}
\vspace*{\fill}
\begin{flushright}
\begin{minipage}{0.55\textwidth}
\raggedleft
\textit{\todo{Inserir epígrafe, se desejado.}}\\[0.3cm]
\todo{Autor da epígrafe, ano}
\end{minipage}
\end{flushright}
\end{epigrafe}


\begin{resumo}
\SingleSpacing
\noindent
\todo{Redigir entre 150 e 500 palavras, em um único bloco, sem citações. Apresentar objetivo, método, principais resultados e conclusões. Atualizar somente quando houver resultados consolidados.}

\vspace{\onelineskip}
\noindent
\textbf{Palavras-chave}: \propellers; otimização geométrica; BEMT; aeroacústica; análise estrutural.
\end{resumo}

\begin{resumo}[ABSTRACT]
\begin{otherlanguage*}{english}
\SingleSpacing
\noindent
\todo{Translate the final Portuguese abstract into English.}

\vspace{\onelineskip}
\noindent
\textbf{Keywords}: propellers; geometric optimization; BEMT; aeroacoustics; structural analysis.
\end{otherlanguage*}
\end{resumo}


\pdfbookmark[0]{\listfigurename}{lof}
\listoffigures*
\cleardoublepage

\pdfbookmark[0]{\listtablename}{lot}
\listoftables*
\cleardoublepage

\begin{siglas}
    \item[BEMT] Blade Element Momentum Theory
    \item[BPF] Blade Passing Frequency
    \item[CFD] Computational Fluid Dynamics
    \item[FEM] Finite Element Method
    \item[LVA] Laboratório de Vibrações e Acústica
    \item[SPL] Sound Pressure Level
    \item[UFSC] Universidade Federal de Santa Catarina
\end{siglas}

\begin{simbolos}
    \item[$B$] Número de pás [-]
    \item[$c$] Corda local da pá [m]
    \item[$D$] Diâmetro do \propeller\ [m]
    \item[$J$] Razão de avanço [-]
    \item[$n$] Velocidade de rotação [rev/s]
    \item[$Q$] Torque [N\,m]
    \item[$R$] Raio do \propeller\ [m]
    \item[$T$] Empuxo [N]
    \item[$\eta$] Eficiência propulsiva [-]
    \item[$\rho$] Massa específica do fluido [kg/m$^3$]
    \item[$\Omega$] Velocidade angular [rad/s]
\end{simbolos}


\pdfbookmark[0]{\contentsname}{toc}
\tableofcontents*
\cleardoublepage

% =============================================================
% ELEMENTOS TEXTUAIS - numeração visível a partir da Introdução.
% =============================================================
\textual
\pagestyle{ufsc}
\aliaspagestyle{chapter}{ufsc}

\chapter{Introdução}
\label{chap:introducao}

\section{Contextualização e problema de pesquisa}
\label{sec:problema}

\todo{Apresentar o contexto de propulsão elétrica, o desafio multidisciplinar do projeto de propellers e formular claramente o problema de pesquisa.}

\section{Justificativa}
\label{sec:justificativa}

\todo{Justificar a relevância de uma ferramenta de projeto preliminar que integre critérios aerodinâmicos, estruturais e aeroacústicos com custo computacional reduzido.}

\section{Objetivos}
\label{sec:objetivos}

\subsection{Objetivo geral}

Desenvolver e validar um algoritmo computacional para otimização geométrica de \propellers, considerando critérios estruturais, aerodinâmicos e aeroacústicos por meio de modelos analíticos implementados computacionalmente.

\subsection{Objetivos específicos}

\begin{itemize}[leftmargin=*]
    \item Identificar os principais parâmetros geométricos, operacionais e materiais que influenciam o desempenho de \propellers;
    \item Desenvolver e consolidar o modelo aerodinâmico baseado em Blade Element Momentum Theory (BEMT) para estimativa de empuxo, torque e eficiência;
    \item Definir e implementar uma abordagem estrutural simplificada para avaliação da integridade da pá;
    \item Consolidar o modelo aeroacústico baseado na formulação de Hanson para estimativa de componentes tonais;
    \item Definir uma estratégia de otimização capaz de maximizar o desempenho aerodinâmico, reduzir a emissão acústica e respeitar restrições estruturais;
    \item Validar os módulos e o algoritmo por comparação com dados experimentais e/ou modelos de referência;
    \item Avaliar as limitações da formulação proposta e indicar possibilidades de continuidade.
\end{itemize}

\section{Estrutura do trabalho}
\label{sec:estrutura_trabalho}

\todo{Ao final da redação, resumir em um parágrafo o conteúdo de cada capítulo.}

\chapter{Revisão bibliográfica}
\label{chap:revisao}

\section{Projeto de propellers}
\label{sec:projeto_propellers}
\todo{Contextualizar o processo de projeto e as principais grandezas de desempenho. O livro de Gur pode servir como referência organizadora desta seção.}

\section{Aerodinâmica de propellers}
\label{sec:aerodinamica}

\subsection{Teoria da quantidade de movimento}
\todo{Apresentar apenas a teoria necessária para construir a BEMT.}

\subsection{Teoria do elemento de pá}
\todo{Definir velocidades, ângulos, forças locais e integração radial.}

\subsection{Blade Element Momentum Theory (BEMT)}
\label{sec:bemt}
\todo{Derivar a formulação adotada, hipóteses, procedimento iterativo e correções utilizadas no código.}

\subsection{Polares aerodinâmicas e XFOIL}
\label{sec:polares}
\todo{Documentar origem/interpolação de $C_L$ e $C_D$, dependência com Reynolds e eventuais limites do XFOIL.}

\section{Modelagem aeroacústica de propellers}
\label{sec:aeroacustica}

\subsection{Ruído tonal e frequência de passagem de pá}
\todo{Apresentar os mecanismos relevantes ao escopo e definir BPF e harmônicos.}

\subsection{Formulação de Hanson}
\label{sec:hanson}
\todo{Apresentar a formulação efetivamente utilizada no código, suas entradas, hipóteses e saídas.}

\section{Modelagem estrutural de pás}
\label{sec:estrutural_teoria}
\todo{Definir o modelo analítico estrutural, esforços considerados, tensões, deslocamentos, torção e critérios de integridade.}

\section{Otimização multidisciplinar de propellers}
\label{sec:otimizacao_teoria}
\todo{Revisar variáveis de projeto, funções objetivo, restrições e métodos de otimização adequados ao problema. A escolha final do algoritmo deve ser justificada mais adiante.}

\section{Síntese da revisão bibliográfica}
\label{sec:sintese_revisao}
\todo{Comparar os trabalhos mais relevantes e explicitar a lacuna que o presente TCC pretende atender. Uma tabela comparativa pode ser útil.}

\chapter{Metodologia}
\label{chap:metodologia}

\section{Visão geral da metodologia}
\label{sec:workflow}
\todo{Inserir o workflow completo: geometria $\rightarrow$ módulos aerodinâmico/estrutural/aeroacústico $\rightarrow$ objetivos e restrições $\rightarrow$ otimizador $\rightarrow$ nova geometria.}

\section{Definição do problema de projeto}
\label{sec:definicao_problema}

\subsection{Parâmetros de entrada e condições operacionais}
\todo{Consolidar limitações geométricas, material, número de pás, dados do motor, velocidade de avanço, rotação e condições atmosféricas.}

\subsection{Variáveis de projeto}
\todo{Definir formalmente os parâmetros efetivamente manipulados pelo otimizador: por exemplo, distribuições de corda e torção e demais variáveis que forem mantidas no escopo final.}

\subsection{Funções objetivo e restrições}
\todo{Formalizar matematicamente objetivos aerodinâmicos/acústicos e restrições estruturais e operacionais.}

\section{Módulo aerodinâmico}
\label{sec:met_bemt}
\todo{Descrever a implementação da BEMT: discretização radial, cálculo iterativo, polares, convergência e outputs $T$, $Q$, potência e eficiência.}

\section{Módulo estrutural}
\label{sec:met_estrutural}
\todo{Descrever o modelo estrutural analítico utilizado dentro do loop de otimização.}

\section{Módulo aeroacústico}
\label{sec:met_hanson}
\todo{Descrever a implementação da formulação de Hanson e o indicador acústico escolhido para a função objetivo.}

\section{Integração dos módulos}
\label{sec:integracao}
\todo{Definir interfaces entre os módulos, unidades, estruturas de dados e sequência de execução.}

\section{Algoritmo de otimização}
\label{sec:met_otimizacao}
\todo{Apresentar o método escolhido, população/pontos iniciais, critérios de parada, tratamento de restrições e reprodutibilidade.}

\section{Estratégia de validação}
\label{sec:validacao}

\subsection{Validação aerodinâmica}
\todo{Comparação com dados experimentais de empuxo, torque e/ou potência.}

\subsection{Validação aeroacústica}
\todo{Comparação com medições de SPL e harmônicos de BPF.}

\subsection{Validação estrutural por FEM}
\todo{Usar o FEM como referência de validação do modelo estrutural simplificado, não necessariamente dentro do loop do otimizador. Incluir análise modal se mantida no escopo final.}

\chapter{Resultados e discussão}
\label{chap:resultados}

% Regra do template UFSC: toda figura/tabela apresentada deve ser discutida no texto.

\section{Validação do módulo aerodinâmico}
\todo{Apresentar casos de validação, métricas de erro e discussão física dos desvios.}

\section{Validação do módulo aeroacústico}
\todo{Comparar resultados tonais com os dados experimentais e discutir tendências e limitações.}

\section{Validação do módulo estrutural}
\todo{Comparar o modelo analítico com FEM e, se aplicável, resultados modais.}

\section{Integração dos módulos}
\todo{Apresentar testes de consistência da arquitetura integrada antes da otimização.}

\section{Estudos paramétricos e otimizações}
\todo{Apresentar os efeitos das variáveis de projeto, frentes de Pareto ou solução selecionada e os compromissos entre desempenho, ruído e integridade estrutural.}

\section{Análise de sensibilidade e consolidação do modelo}
\todo{Quantificar a sensibilidade das respostas e discutir robustez, custo computacional e limitações.}

\chapter{Conclusões}
\label{chap:conclusoes}

\todo{Responder diretamente ao objetivo geral e aos objetivos específicos, destacando o que foi validado, os principais resultados e as limitações observadas.}

\section{Sugestões para trabalhos futuros}
\todo{Registrar extensões que ficaram fora do escopo: modelos de maior fidelidade, refinamentos estruturais/aeroacústicos, novas estratégias de otimização e campanhas experimentais adicionais.}


% =============================================================
% ELEMENTOS PÓS-TEXTUAIS
% =============================================================
\postextual

\bibliography{references}

\begin{apendicesenv}

% O template UFSC apresenta explicitamente "APÊNDICE A - ...".
\renewcommand{\printchapternum}{APÊNDICE~\thechapter\space-\space}
\chapter{Informações complementares do algoritmo}
\label{ap:algoritmo}

\todo{Usar este apêndice para conteúdo produzido por você que seja importante para reprodutibilidade, mas excessivamente detalhado para o corpo principal: pseudocódigos, parâmetros completos, testes adicionais etc.}

\end{apendicesenv}

% \begin{anexosenv}

\renewcommand{\printchapternum}{ANEXO~\thechapter\space-\space}
\chapter{Documento complementar}
\todo{Anexos contêm documentos não elaborados pelo autor que sejam necessários para a compreensão do trabalho.}

\end{anexosenv}


\end{document}
